(1) 未扰动哈密顿量 $H=v^2/2-\omega_0^2x^2/2+\beta x^4/4$。令 $H=0$，分离变量并取同宿支，得到 $x_h=\sqrt2\omega_0\operatorname{sech}(\omega_0t)/\sqrt\beta$、$v_h=-\sqrt2\omega_0^2\operatorname{sech}(\omega_0t)\tanh(\omega_0t)/\sqrt\beta$。代回方程可直接验证；当 $t\to\pm\infty$ 时二者归零。

(2) 沿受扰运动求 $\dot H=v(\ddot x-\omega_0^2x+\beta x^3)=\epsilon v[f\cos(\Omega t)-\delta v]$。令扰动相位为 $t_0$，一阶轨道能量差 $M(t_0)=f\int_{-\infty}^{\infty}v_h(t)\cos[\Omega(t+t_0)]dt-\delta\int_{-\infty}^{\infty}v_h^2(t)dt$，简单零点预示横截交叉。

(3) 置 $u=\tanh\omega_0t$，得 $\int v_h^2dt=4\omega_0^3/(3\beta)$。由 $\int_{-\infty}^{\infty}\operatorname{sech}(\omega_0t)e^{i\Omega t}dt=(\pi/\omega_0)\operatorname{sech}(\pi\Omega/(2\omega_0))$ 并对 $v_h=\dot x_h$ 分部积分，得到驱动积分 $\pi\sqrt2\Omega\sin(\Omega t_0)\operatorname{sech}(\pi\Omega/(2\omega_0))/\sqrt\beta$。

(4) 因而 $M(t_0)=A\sin\Omega t_0-D$，其中 $A=\pi\sqrt2 f\Omega\operatorname{sech}(\pi\Omega/(2\omega_0))/\sqrt\beta$，$D=4\delta\omega_0^3/(3\beta)$。若 $A>D$ 存在简单零点，足够小的 $\epsilon$ 下预测稳定和不稳定流形横截相交。$A=D$ 是切触边界，一阶理论不能据此判定混沌。
